\section{YÊU CẦU VÀ MỤC TIÊU ĐỀ TÀI}
\subsection{Yêu cầu đề tài}
Đề tài yêu cầu nghiên cứu và triển khai một hệ thống chăm sóc sức khỏe điện tử tích hợp với thiết bị cân thông minh như \textbf{Xiaomi Smart Scale 2} và \textbf{CRENOT Gofit S2}. Thiết bị này sẽ cung cấp các chỉ số sức khỏe cơ bản như cân nặng, tỷ lệ mỡ cơ thể, khối lượng cơ, nước, và xương thông qua công nghệ phân tích trở kháng điện sinh học (BIA). Hệ thống cần đảm bảo các yêu cầu sau:

\begin{enumerate}
    \item \textbf{Thu thập và quản lý dữ liệu tự động:} Hệ thống phải đảm bảo khả năng thu thập dữ liệu từ thiết bị cân sức khoẻ thông minh và lưu trữ chúng trong 1 file csv với theo từng cá nhân riêng biệt.
    \item \textbf{Vẽ biểu đồ, trực quan hoá dữ liệu:} Dữ liệu được thu thập sẽ được lưu lại trong 1 file csv với mục đích dùng để vẽ các biểu đồ về thông số sức khoẻ qua các ngày khác nhau, với mục đích để tiện cho việc theo dõi của người dùng
    \item \textbf{Phân tích dữ liệu và đưa ra khuyến nghị:} Dựa trên các dữ liệu thu thập, hệ thống cần có khả năng phân tích, đưa ra các khuyến nghị về chế độ dinh dưỡng, tập luyện, và cảnh báo sớm về các vấn đề sức khỏe cho người dùng.
    \item \textbf{Bảo mật và bảo vệ dữ liệu cá nhân:} Hệ thống cần tuân thủ các tiêu chuẩn bảo mật quốc tế về bảo vệ thông tin cá nhân, đảm bảo an toàn cho dữ liệu sức khỏe của người dùng trong quá trình truyền tải và lưu trữ.
\end{enumerate}

\subsection{Mục tiêu đề tài}
Mục tiêu của đề tài là xây dựng một hệ thống chăm sóc sức khỏe điện tử thông minh, giúp người dùng theo dõi và quản lý sức khỏe cá nhân một cách hiệu quả và tiện lợi. Cụ thể:

\begin{enumerate}
    \item \textbf{Tạo nền tảng quản lý sức khỏe toàn diện:} Hệ thống giúp người dùng theo dõi các chỉ số sức khỏe quan trọng từ thiết bị cân sức khoẻ thông minh, đồng thời cung cấp phân tích chi tiết về tình trạng sức khỏe của họ.
    \item \textbf{Cải thiện chất lượng cuộc sống:} Thông qua việc cung cấp thông tin và khuyến nghị sức khỏe, hệ thống sẽ giúp người dùng duy trì lối sống lành mạnh và ngăn ngừa các bệnh liên quan đến cân nặng, chế độ ăn uống và tập luyện không hợp lý.
    \item \textbf{Phát triển giải pháp y tế số:} Đề tài sẽ đóng góp vào sự phát triển của các giải pháp y tế số tại Việt Nam, giúp các cơ sở y tế và bệnh viện quản lý sức khỏe từ xa cho bệnh nhân một cách hiệu quả hơn, đặc biệt trong các trường hợp bệnh mãn tính.
    \item \textbf{Tăng cường nhận thức về chăm sóc sức khỏe:} Hệ thống sẽ giúp nâng cao nhận thức về việc theo dõi và duy trì sức khỏe hàng ngày của cá nhân, khuyến khích họ thường xuyên kiểm tra các chỉ số quan trọng để phát hiện sớm các dấu hiệu bất thường.
\end{enumerate}

Source code của bài báo cáo sẽ được lưu trữ và cập nhật tại đây: \href{https://github.com/DatTwenty3/xiaomi_smart_scale_2}{DatTwenty3 - GitHub}